%%%%%%%%%%%%%%%%%%%%%%%%%%%%%%%%%%%%%%%%%%%%%%%%%%%%%%%%%%%%%%%%%%%%%%%%%%%
%%  preliminaries.tex  –  Титульный лист, Аннотация, Благодарности
%%  v10 – Обновлено для нового названия диссертации.
%%        Аннотация полностью переработана в стиле концепций/теорий.
%%        Никаких \cite{} в этом файле для корректной нумерации [1].
%%%%%%%%%%%%%%%%%%%%%%%%%%%%%%%%%%%%%%%%%%%%%%%%%%%%%%%%%%%%%%%%%%%%%%%%%%%

%% ====================== ТИТУЛЬНЫЙ ЛИСТ ======================================
\begin{titlepage}
\thispagestyle{empty}
\begin{center}

{\small
МИНИСТЕРСТВО НАУКИ И ВЫСШЕГО ОБРАЗОВАНИЯ\\
РОССИЙСКОЙ ФЕДЕРАЦИИ\\
ФЕДЕРАЛЬНОЕ ГОСУДАРСТВЕННОЕ АВТОНОМНОЕ ОБРАЗОВАТЕЛЬНОЕ\\
УЧРЕЖДЕНИЕ ВЫСШЕГО ОБРАЗОВАНИЯ\\[2pt]
\textbf{<<Национальный исследовательский ядерный университет <<МИФИ>>}\\
\textbf{(НИЯУ МИФИ)}
}

\vspace{1.8cm}

{\small
Отчёт <<Научно-исследовательская деятельность аспиранта и подготовка\\
к защите диссертации на соискание учёной степени кандидата наук>> за
}

\vspace{0.5cm}

{\small
\underline{\hspace{1.2cm}}\textbf{6}\underline{\hspace{1.2cm}} семестр
}

\vspace{1.5cm}

{\small\textbf{<<Разработка состязательных устойчивых моделей в гибридных системах обнаружения и предотвращения вторжений>>}}

\vspace{2.5cm}

\begin{minipage}{0.82\textwidth}
\small\noindent
\begin{tabular}{@{}p{0.36\textwidth} p{0.60\textwidth}@{}}
Аспирант & \textbf{Анаэдевха Роджер Ник}\\
Научная специальность & 2.3.1 Системный анализ, управление и\\
 & обработка информации, статистика\\[1.2cm]
Научный руководитель & \textbf{Трофимов Александр Геннадьевич}\\
Должность, степень, звание & Доктор физико-математических\\
 & наук, профессор\\[1.2cm]
Дата защиты: & \\[0.6cm]
Результат защиты: & \\
\end{tabular}
\end{minipage}

\vfill

{\small Москва, 2026}
\end{center}
\end{titlepage}


%% ====================== АННОТАЦИЯ ========================================
\chapter*{Аннотация}
\addcontentsline{toc}{chapter}{Аннотация}

Настоящая диссертация посвящена разработке состязательных устойчивых моделей в гибридных системах обнаружения и предотвращения вторжений (IDPS). Модели графовых нейронных сетей (ГНС) являются основной парадигмой глубокого обучения для извлечения предсказательных представлений из реляционных данных, однако их робастность в условиях целенаправленных атак отстаёт от точности предсказаний на 25--65 процентных пунктов в критически важных системах: здравоохранении, финансах, промышленном IoT, автономных системах и кибербезопасности. Диссертация решает эту фундаментальную проблему через разработку семи новых методов, действующих при трёх рабочих условиях --- состязательная среда, нестационарные условия и распределённая сетевая динамика --- при одновременном выполнении четырёх сквозных требований: робастности, точности, эффективности и дифференциальной приватности.

Впервые установлен принцип сертифицированной устойчивости моделей ГНС непрерывного времени (CT-TGNN и SDE-TGNN), основанный на том, что динамика представлений узлов графа, управляемая дифференциальными уравнениями с ограниченной гладкостью, порождает доказуемые границы чувствительности выходных сигналов к возмущениям входных данных и структуры. Стохастическое расширение SDE-TGNN заменяет детерминированную динамику стохастическими дифференциальными уравнениями Ито и выводит аналитическое распространение моментов из уравнения Фоккера--Планка, достигая средневзвешенного F1 99,0\% на трёх бенчмарках общим объёмом более 52 миллионов образцов, ожидаемой ошибки калибровки 0,042, сертифицированной состязательной робастности 76,2\% при радиусе возмущения 0,25 и кросс-доменного переноса 90,5\% F1 без дообучения.

Предложена оригинальная теория совместного обеспечения византийской устойчивости, формальной приватности и распределительного выравнивания (FedLLM-API) при обучении моделей ГНС на разрозненных фрагментах графов, принадлежащих разным организациям, без передачи исходных данных. Сходимость доказана со скоростью $\calO(1/\sqrt{T})$ при 30\% злонамеренных участников, а интеграция большой языковой модели обеспечивает 87,6\% F1 при обнаружении новых категорий угроз в режиме zero-shot.

Разработана новая концепция многоуровневого графового представления (TripleE-TGNN), в которой модели ГНС одновременно обучают структуру на уровнях сервисов, трассировок и узлов с механизмом защиты от утраты ранее усвоенных закономерностей через упругую консолидацию весов, достигая точности 96,8\% при снижении вычислительных затрат на 40\%.

Предложен новый подход к эффективному выводу (MambaShield) на основе селективных пространств состояний, снижающий вычислительную сложность с квадратичной до почти линейной через свёртки быстрого преобразования Фурье, при сохранении устойчивости к отравлению обучающих данных: точность 91,4\% при 25\% отравления с PAC-байесовскими сертификатами обобщения и задержкой вывода 0,5~мс на поток.

Впервые предложена интегрированная теория калиброванной неопределённости для моделей ГНС (стохастический трансформер и UC-HGP), позволяющая разделять источники ошибок на устранимые (эпистемические) и неустранимые (алеаторные) с ожидаемой ошибкой калибровки 0,078, обеспечивая снижение нагрузки на аналитиков на 43\% через автоматическую сортировку оповещений по неопределённости.

Построена оригинальная концепция единой сертификации устойчивости, в которой взаимодействие защитника и противника моделируется как стратегическая игра через равновесия Штакельберга с онлайн-обучением без регрета, распространение интервальных границ и рандомизированное сглаживание. Система связывает все семь методов через регуляризацию согласованности с доказанной скоростью сходимости.

Разработана новая теория предметно-ориентированного фундаментального моделирования (CyberSecLLM) на графовых структурах с блоками селективного пространства состояний, дополненными ОДУ, достигающая средней точности 89,1\% в режиме zero-shot на бенчмарке из 11 задач (на 20,5 п.п. выше GPT-4) при обработке графовых последовательностей длиной в миллион токенов за 2,3 секунды.

Полная система реализована как платформа промышленного уровня RobustIDPS.ai (DOI: 10.5281/zenodo.19129512), включающая 46 страниц, 9 функциональных групп и 13+ интегрированных моделей нейронных сетей. Четырёхуровневая система защиты LLM достигает 94,5\% макроусреднённого обнаружения по 517 сценариям атак с четырьмя коммерческими бэкендами (Claude, GPT-4o, Gemini, DeepSeek). Модуль Multi-Agent PQC-IDS обеспечивает классификацию постквантового трафика. Платформа развёртывается как плагин для Snort и Suricata.

Комплексная экспериментальная валидация охватывает шесть бенчмарк-наборов данных общим объёмом 84,2 миллиона размеченных образцов по 84 категориям атак, единую систему оценки из 23 метрик, абляционные исследования, оценки состязательной робастности при шести семействах атак и эксперименты по кросс-доменному переносу на речь и здравоохранение. Единая система превосходит лучшие базовые методы на 2,9--5,6 процентных пунктов на всех наборах данных.

Все методы сформулированы предметно-независимо на уровне математических операций над графовыми структурами и валидированы через применение к гибридным системам обнаружения и предотвращения вторжений. Все 12 пересечений матрицы условие--требование покрыты с подтверждённой производительностью. Акты внедрения получены от Positive Technologies, Google Cloud Platform и Suricata Open Source IDPS.

Диссертация состоит из введения, пяти глав, заключения, библиографического списка из 190 наименований и трёх приложений. Объём основного текста составляет 210 страниц.


%% ====================== БЛАГОДАРНОСТИ ================================
\chapter*{Благодарности}
\addcontentsline{toc}{chapter}{Благодарности}

Выражаю искреннюю благодарность моему научному руководителю, профессору Александру Геннадьевичу Трофимову, чьё руководство, интеллектуальная строгость и щедрая поддержка сопровождали каждый этап данного исследования. Его глубокая экспертиза в области математического моделирования и системного анализа заложила основу, на которой были построены теоретические вклады настоящей диссертации.

Благодарю профессорско-преподавательский состав и коллег-исследователей Института интеллектуальных кибернетических систем и кафедры кибернетики Национального исследовательского ядерного университета <<МИФИ>> за создание стимулирующей академической среды. Обсуждения, семинары и дух сотрудничества на кафедре существенно обогатили данную работу.

С благодарностью отмечаю поддержку аспирантской программы МИФИ, гранта на создание исследовательских центров в области искусственного интеллекта от Министерства экономического развития Российской Федерации (идентификатор 000000C313925P3Q0002), а также вычислительные ресурсы, предоставленные для проведения экспериментальной работы, описанной в диссертации. Масштаб эмпирической валидации был бы невозможен без доступа к соответствующей инфраструктуре.

Наконец, я бесконечно признателен моей семье и друзьям, чьё терпение, поддержка и вера в это начинание поддерживали меня на протяжении всех лет учёбы и исследований. Спасибо всем.

\hfill Роджер Ник Анаэдевха

\hfill Москва, 2026 

